%!TEX root = ../main.tex

\chapter{Introduction\label{chap:introduction}}

% First, introduce the reader to the research topic.
% Start with the most general view and slowly converge to the particular field, sub-field, and the challenges you face.
% You can cite others' work here \cite{baca2021mrs}.

In the context of drone racing, where quadcopters are pushed to their physical and computational limits, there is a growing trend toward using neural networks for control.
As demonstrated in \cite{baca2021mrs}, this approach has made it possible to outperform some of the world's most skilled human pilots.

These neural networks controllers are tipically trained using \ac{RL} algorithms.
Due to the sample ineffiency of \ac{RL} and the high cost of collecting real-world flight data, training is performed in simulated environments that provide a safe and scalable alternative.
Despite this, the resulting control policies often perform worse than expected or fail to transfer to real-world scenarios altogether.
This discrepancy in the policy's performance between simulation and reality is known as the sim-to-real gap, and it mainly arises from an undesired overfit of the learned policy to the simulated environment.

% To mitigate this issue, improving the system's identification, using digital twins, or randomizing the domain, among other strategies, have shown to significantly reduce the sim-to-real gap.

% In this semesteral project, the focus has been placed on understanding and implementing \ac{DR} in learning-based control policies for quadcopter racing competitions.
% Key parameters from the drone's dynamics, perception, and observation spaces have been identified and randomized under different \ac{DR} configurations.
% Finally, {\color{red}{X}} randomized policies are analyzed and compared to non-randomized ones. Their performance is evaluated to assess how DR improves adaptability and transferability to real-world scenarios.

\section{Related works}

This section should contain related state-of-the-art works and their relation to the author's work.
We usually cite the original works like this \cite{benallegue2008high}.
You can also cite multiple papers at once like this \cite{baca2016embedded, baca2021mrs}.

Demonstrating that Domain Randomization (DR) can significantly increase the adaptability and robustness of autonomous systems, the work of Ferede et al. \cite{ferede2025} serves as the direct inspiration for this project. 
Their research showed that a DR-trained neural network controller could successfully generalize across various quadcopter platforms.
\section{Contributions}

(Note: rather than contribute this work aims to reproduce the results from \cite{ferede2025})

\section{Mathematical notation}

It is a good practice to define basic mathematical notation in the introduction.
See \reftab{tab:mathematical_notation} for an example.

\begin{table*}[!h]
  \scriptsize
  \centering
  \noindent\rule{\textwidth}{0.5pt}
  \begin{tabular}{lll}
    $\mathbf{x}$, $\bm{\alpha}$ & vector, pseudo-vector, or tuple\\
    $\mathbf{\hat{x}}$, $\bm{\hat{\omega}}$& unit vector or unit pseudo-vector\\
    $\mathbf{\hat{e}}_1, \mathbf{\hat{e}}_2, \mathbf{\hat{e}}_3$ & elements of the \emph{standard basis} \\
    $\mathbf{X}, \bm{\Omega}$ & matrix \\
    $\mathbf{I}$ & identity matrix \\
    $x = \mathbf{a}^\intercal\mathbf{b}$ & inner product of $\mathbf{a}$, $\mathbf{b}$ $\in \mathbb{R}^3$\\
    $\mathbf{x} = \mathbf{a}\times\mathbf{b}$ & cross product of $\mathbf{a}$, $\mathbf{b}$ $\in \mathbb{R}^3$\\
    $\mathbf{x} = \mathbf{a}\circ\mathbf{b}$ & element-wise product of $\mathbf{a}$, $\mathbf{b}$ $\in \mathbb{R}^3$ \\
    $\mathbf{x}_{(n)}$ = $\mathbf{x}^\intercal\mathbf{\hat{e}}_n$ & $\mathrm{n}^{\mathrm{th}}$ vector element (row), $\mathbf{x}, \mathbf{e} \in \mathbb{R}^3$\\
    $\mathbf{X}_{(a,b)}$ & matrix element, (row, column)\\
    $x_{d}$ & $x_d$ is \emph{desired}, a reference\\
    $\dot{x}, \ddot{x}, \dot{\ddot{x}}$, $\ddot{\ddot{x}}$ & ${1^{\mathrm{st}}}$, ${2^{\mathrm{nd}}}$, ${3^{\mathrm{rd}}}$, and ${4^{\mathrm{th}}}$ time derivative of $x$\\
    $x_{[n]}$ & $x$ at the sample $n$ \\
    $\mathbf{A}, \mathbf{B}, \mathbf{x}$ & LTI system matrix, input matrix and input vector\\
    \emph{SO(3)} & 3D special orthogonal group of rotations\\
    \emph{SE(3)} & \emph{SO(3)}~$\times~\mathbb{R}^3$, special Euclidean group\\
  \end{tabular}
  \noindent\rule{\textwidth}{0.5pt}
  \caption{Mathematical notation, nomenclature and notable symbols.}
  \label{tab:mathematical_notation}
\end{table*}
